\chapter{Foley Catheter Insertion}

\section*{Overview}
A Foley catheter is a flexible tube inserted through the urethra into the bladder to continuously drain urine. It is held in place by an inflated balloon at the tip. The exact approach, setting, and preparation depend on the indication, anatomy, and the patient's health.

\section*{Alternative Names}
Urinary catheterization, Indwelling catheter insertion

\section*{Principle}
A sterile catheter with a deflated balloon at its tip is passed through the urethra into the bladder. The balloon is then inflated with sterile water to hold the catheter in place, providing a closed drainage pathway for urine.
\section*{History}
Catheters have been used since antiquity. Frederic Foley designed the modern balloon retention catheter in 1929, which was commercialized in the late 1930s.

\section*{Why It Is Done}
Ordered for acute urinary retention, monitoring urine output in critically ill patients, during surgical procedures, or for bladder decompression.

\section*{Is This a Primary or Secondary Test or Procedure}
Primary treatment for acute urinary retention and for precise urine output monitoring.

\section*{How It Is Performed}
Under sterile conditions. The periurethral area is cleaned, a sterile lubricant is applied, and the catheter is gently inserted through the urethra into the bladder. The balloon is inflated with sterile water.

\section*{Preparation}
Ensure sterile gloves, catheter kit, and antiseptic prep are available. Position the patient supine (female in dorsal recumbent).

\section*{Contraindications and Precautions}
Suspected urethral disruption (typically in pelvic trauma; look for blood at the meatus or high-riding prostate).

\section*{Risks}
Urinary tract infection (CAUTI), urethral trauma/bleeding, false passage creation, or bladder spasms.

\section*{Aftercare and Recovery}
Maintain a closed drainage system, secure the catheter to the patient's thigh, and keep the drainage bag below the level of the bladder.

\section*{Outcome and Follow-up}
Immediate flow of clear urine confirms correct placement. Monitor for patency and urine volume.

\section*{Expected Results}
Continuous drainage of clear yellow urine; no leakage around the catheter; no patient distress.

\section*{Follow-up}
Regular catheter hygiene, monitoring for UTI symptoms, and prompt removal when no longer indicated.

\section*{When to Seek Medical Care}
Follow the procedure team's specific instructions. Seek urgent medical assessment for severe or worsening pain, uncontrolled bleeding, fever or signs of infection, trouble breathing, chest pain, fainting, new weakness or confusion, inability to urinate, or any rapidly worsening symptom. The appropriate warning signs depend on the procedure and the patient's medical condition.

\section*{References}
\begin{enumerate}
  \item MedlinePlus. Foley Catheter Insertion. U.S. National Library of Medicine; 2025.
  \item Current Medical Diagnosis and Treatment. McGraw-Hill; 2025.
\end{enumerate}

\clearpage
